This chapter studies how computation behaves when gates, wires, or intermediate states are subject to perturbation. The central quantity is a reliability function \(\varepsilon(r)\), interpreted as the probability of correct operation as a function of a resource parameter \(r\), such as depth, size, code distance, or redundancy level. In the simplest form, one may write
\[
\varepsilon(r) = 1 - \Pr[\text{error}].
\]
The point of the chapter is not merely to measure failure, but to understand how compositional structure can be made robust under noise, and when redundancy can be organized so that reliability improves rather than degrades as systems scale.

A noisy computation model treats elementary operations as imperfect. A gate may flip its intended output with some small probability, a wire may corrupt a signal, and a composed circuit may accumulate error across layers. In this setting, the design problem becomes one of robust closure: given a family of operations closed under composition in the idealized setting, how can one construct a corresponding family of fault-tolerant implementations whose effective behavior remains within the same computational class? The answer is typically not to eliminate noise entirely, but to encode information and computation so that local faults are detected, corrected, or rendered harmless by aggregation.

A classical strategy is redundancy. If a bit is represented by several physical bits, then a decoding rule can recover the intended logical value even when some fraction of the physical bits are corrupted. The simplest example is majority logic: if three copies of a bit are transmitted or stored, then the majority vote returns the correct value whenever at most one copy is faulty. More generally, majority-based schemes can be iterated hierarchically, yielding recursive constructions in which each level of encoding suppresses the effective error rate of the level below. Such schemes are closely related to von Neumann-style reliable computation, in which unreliable components are assembled into a more reliable whole by combining redundancy with local correction.

Von Neumann schemes are especially important because they show that reliable computation is not limited to perfectly reliable primitives. Instead, one can build a threshold phenomenon: if the physical error rate is below a certain bound, then sufficiently deep or sufficiently redundant constructions can drive the logical error rate arbitrarily low. This threshold behavior is one of the chapter's central themes. It provides a quantitative closure principle: below threshold, the class of computable functions is stable under noisy implementation, in the sense that the intended logical operation can be recovered with high probability after encoding, composition, and decoding.

Error-correcting codes provide a more systematic language for this phenomenon. A code is not only a storage device but also a compositional symmetry: it identifies a structured subset of states together with encoding and decoding maps that preserve logical information under a controlled family of perturbations. From this perspective, the code space is an invariant substructure, and the correction procedure is a projection back onto that invariant. The algebraic viewpoint is useful because it emphasizes that fault tolerance is not an ad hoc patch; it is a design principle based on symmetry, invariance, and controlled quotienting of noisy degrees of freedom.

This perspective also clarifies why codes compose. If one encoding protects against a certain class of local errors and a second encoding protects the encoded representation against higher-level faults, then the combined system can inherit robustness across multiple scales. Such nested constructions are the reliability analogue of hierarchical modularity: each layer abstracts away the details of the layer below while preserving the logical content. In practice, this means that the same conceptual tools used earlier in the book to study compositional closure can be applied to fault tolerance, with the additional requirement that the relevant equivalences be probabilistic rather than exact.

Threshold theorems formalize this intuition. They assert, under suitable assumptions on the noise model and the available correction gadgets, that there exists a constant \(p_{\ast}\) such that if the physical error rate \(p < p_{\ast}\), then arbitrarily long computations can be performed with arbitrarily small logical error by increasing redundancy and using recursive correction. The exact value of the threshold depends on the architecture, the noise model, and the correction protocol, but the qualitative conclusion is robust: reliable computation is possible in principle even when every component is imperfect, provided the imperfections are sufficiently sparse and sufficiently well-behaved.

The chapter is also concerned with how reliability scales with resource usage. The function \(\varepsilon(r)\) is intended to capture this dependence, whether \(r\) denotes circuit depth, code distance, number of redundant copies, or another measure of implementation scale. In a favorable regime, \(\varepsilon(r)\) increases toward one as \(r\) grows, reflecting improved reliability. In an unfavorable regime, additional depth may simply accumulate more opportunities for failure. The design challenge is therefore to choose encodings and correction mechanisms so that the gain from redundancy dominates the cost of extra operations.

A practical artifact for this chapter is a fault-injection harness. Such a harness perturbs a circuit or abstract computation according to a specified noise model, runs repeated trials, and records empirical error rates as a function of redundancy parameters. The resulting error-versus-redundancy plots provide an operational view of threshold behavior: one can observe whether additional copies, larger code distance, or deeper correction layers reduce the logical error rate, and at what point diminishing returns set in. These plots are not merely illustrative; they are part of the chapter's reproducibility protocol and serve as a bridge between theoretical claims and empirical validation.

Taken together, noisy gates, majority logic, error-correcting codes, and threshold theorems show that closure in digital systems must sometimes be understood probabilistically. The ideal algebra of composition remains relevant, but it is supplemented by a robustness calculus that tracks how errors propagate, how they are suppressed, and when they can be made asymptotically negligible. This chapter develops that viewpoint as a foundation for later discussions of reliable architectures, fault-tolerant program semantics, and symmetry-aware design under uncertainty.
